%!TEX program = lualatex
\documentclass[ngerman,headheight=70pt]{scrartcl}
\usepackage[ngerman]{babel}
\usepackage{blindtext}
\usepackage[utf8]{luainputenc}
\usepackage{scrlayer-scrpage}
\usepackage[T1]{fontenc}
\usepackage{fontspec}
\usepackage{multicol}
\usepackage{enumitem}
\usepackage{epigraph}
\usepackage[german=quotes]{csquotes}
\usepackage[bottom=2cm,footskip=8mm]{geometry}
\setmainfont{Source Sans Pro}
\MakeOuterQuote{"}

\setlength{\parindent}{0pt}
\setlength{\parskip}{6pt}

\pagestyle{plain.scrheadings}
\thispagestyle{scrheadings}
\lohead{\vspace*{-1.5cm}\includegraphics[scale=0.4]{../images/stupa_logo_rot.pdf}}
\rohead{\\Präsidium des\\Studierendenparlaments}

\lofoot[Universität Hamburg · Präsidium des Studierendenparlaments\\\hrule Von-Melle-Park 5 · D-20146 Hamburg]{Universität Hamburg · Präsidium des Studierendenparlaments\\\hrule Von-Melle-Park 5 · D-20146 Hamburg}
\cofoot[]{}
\cohead[\thepage]{}

\renewcommand{\thesubsection}{TOP \arabic{subsection}}
\renewcommand{\thesubsubsection}{\arabic{subsubsection}.}
\setcounter{subsection}{-1}

\begin{document}
    UHH · StuPa-Präsidium · Von-Melle-Park 5 · D-20146 Hamburg

    \section*{Beschlüsse des Studierendenparlamentes vom 10. November 2016}

    Es gibt insgesamt 2 beschlossene Texte zum Thema bedarfsgerechte
    Hochschulfinanzierung. Der erste mit dem Titel "Für eine bedarfsgerechte und
    zukunftsfähige Hochschulfinanzierung" wurde mit 7 Dafürstimmen, 0 Gegenstimmen
    und mehrheitlichen Enthaltungen angenommen.
    Der zweite Text mit dem Titel "Für eine bedarfsgerechte und nachhaltige
    Hochschulfinanzierung" wurde mit 11 Dafürstimmen, einer Gegenstimme und
    vielen Enthaltungen angenommen.

    \subsection*{Für eine bedarfsgerechte und zukunftsfähige Hochschulfinanzierung}

    Im Dezember dieses Jahres verabschiedet die Hamburgische Bürgerschaft den
    Doppelhaushalt für die Jahre 2017/18. Im letzten Doppelhaushalt 2015/16 bekam
    die Universität 276.956 Tsd. Euro im Jahr 2015 und 279.393 Tsd. Euro im Jahr
    2016. Diese Steigerung entsprach genau 0,88%.

    Eine Steigerung von 0,88\% war und ist nicht einmal genug, um die Lohn- und
    Gehaltsteigerungen sowie die allgemeine Teuerungsrate auszugleichen. Es
    entspricht somit einer realen Kürzung, auch wenn der nominelle Beitrag
    ansteigt. Vor dem Hintergrund unserer Forderung nach einer bedarfsgerechten
    Finanzierung ist so eine Steigerung daher in jedem Fall unzureichend. Der
    Wissenschaftsrat hat dies ebenfalls erkannt und im Anschluss an die Begutachtung
    des MINT-Bereichs einen jährlichen Anstieg des Budgets von 3,5\% gefordert.
    Zwar erhöht sich der Betrag für die Universität im Vergleich zum Jahr 2016 um
    3,5\%, aber die Uni muss im Gegenzug auch gewisse Abschreibungen neu im
    Haushalt ausweisen. Defacto entspricht die Erhöhung daher weiterhin nur 0,88\%
    und ist somit weiterhin nicht annähernd genügend.

    Ein Anstieg von 3,5\% wäre zwar ein erster Schritt in die richtige Richtung,
    reicht aber natürlich noch nicht aus. Wir fordern weiterhin eine
    bedarfsgerechte Finanzierung der Universität Hamburg, die nicht zulasten
    anderer wissenschaftlicher oder sozialer Einrichtungen geht.

    Die nächste Bürgerschaftswahl ist erst 2020. Es wird also vorher noch ein
    Doppelhaushalt (2019/20) verabschiedet werden. Dort fordern wir mindestens
    eine Erhöhung um 3,5\%, um der Gesamtforderung einer bedarfsgerechten
    Finanzierung zumindest schrittweise näher zu kommen.

    Perspektivisch ist jedoch alleinig die bedarfsgerechte Ausfinanzierung der
    Universität und weiterer wissenschaftlicher Einrichtungen das anzustrebende
    Ziel. Die Zukunft hängt von der Bildung ab und gerade angesichts des
    Klimawandels und dessen Einfluss auf die Schifffahrt wird der Hamburger
    Hafen nicht auf Dauer die Lösung für Hamburg sein. Auch einzelne
    Leuchtturmprojekte oder Exzellenzcluster werden dann nicht helfen.
    Hamburg ist bereits jetzt eine der beliebtesten Städte unter Studierenden.
    Dieses Jahr haben über 6000 neue Studierende an der Universität Hamburg
    angefangen. Diese Zahl wird sich eher nur noch steigern. Spitzenförderung
    reicht also nicht aus. Es bedarf schlichtweg einer Förderung in der Breite
    und die ist nur über eine Ausfinanzierung der Universität im Rahmen des
    Globalbudgets zu erreichen.

    Diese Vielzahl an Studierenden bietet ein großes Potential, wenn forschendes
    Lehren und Lernen in der Breite umgesetzt würde. Denn es gibt genügend
    gesellschaftlich relevante Probleme, die produktiv an der Universität und
    den Hochschulen im Rahmen des Studiums bearbeitet werden könnten.
    Krebsforschung, Datenschutz und Klimawandel seien nur einmal exemplarisch
    genannt. Damit dies jedoch gelingen kann und ein wirkliches interdisziplinäres
    Forschen, Lehren und Lernen in einer Einheit gelingen kann, bedarf es einer
    flächendeckend ausreichenden Finanzierung.

    Hamburg ist die Stadt mit den meisten Millionären Deutschlands und ein "Tor
    zur Welt". Es sollte daher möglich sein die Universität Hamburg und die
    weiteren Hamburger Hochschulen und Universitäten entsprechend dieses Ziels
    ausreichend zu finanzieren. Wer sich dies nicht klar auf die Fahnen schreibt,
    verspielt die Zukunft der nächsten Generation.

    Die Zeit des Sparens in diesem Bereich muss vorbei sein. Wissenschaft ist
    kein Schönwetterthema, welches geleistet werden kann, wenn es sonst nichts
    wichtiges gibt. Es ist ein elementar wichtiger Bereich für die Stadt Hamburg,
    Deutschland und die Welt. Wenn Hamburg ein "Wissenschaftsstandort" werden
    soll, dann muss dieser Bereich die gleiche Priorität und den gleichen
    Stellenwert bekommen, wie ihn momentan der Hafen bzw. die Wirtschaft
    allgemein hat.

    \newpage

    \subsection*{Für eine bedarfsgerechte und nachhaltige Hochschulfinanzierung}

    Im Dezember dieses Jahres wird die Hamburgische Bürgerschaft den
    Doppelhaushalt für die Jahre 2017/18 verabschieden. Der aktuelle Entwurf
    fügt sich dabei lückenlos in die bisherige Linie des Hamburgischen Senats
    ein, die öffentlichen Bereiche mit Steigerungen um ungefähr 0.88\% abzuspeisen.

    Das war und ist nicht einmal ausreichend, um die Lohn- und Gehaltssteigerungen
    sowie die allgemein Inflation auszugeichen. Die vermeintlichen Steigerungen
    entsprechen somit einer realen Kürzung! Schauen wir exemplarisch auf die
    Situation der Hochschulen, so wird schnell klar, dass sich im Rathaus wider
    besseres Wissen und wachsenden Protest dem Dogma der "schwarzen Null"
    unterworfen wird.

    Die schon heute dramatische Unterfinanzierung der Wissenschaftslandschaft
    führt ungebremst in die falsche Richtung: Sinkende Betreuungsrelationen,
    steigende Drittmittelabhängigkeit und prekäre Beschäftigungsbedingungen sowie
    die durch Unterfinanzierung forcierte Konkurrenz behindern die
    Hochschulmitglieder darin, ihrer Aufgabe wissenschaftlicher Reflexion und
    kritischer Intervention nachzukommen. Die Studienbedingungen der Bologna-Reform,
    darauf ausgerichtet Akademiker*innen als Humankapital "am Fließband" zu
    produzieren, überzeugen niemanden mehr und schmälern das Potential der
    Hochschulen, ein Ort der Bildung demokratischer Persönlichkeiten zu sein.
    Dagegen gibt es inzwischen weit verbreitete und entwickelte Kritik: Mit der
    Schuldenbremse soll der Eindruck erweckt werden, es fehle an den Mitteln,
    etwas Besseres zu verwirklichen.

    Der Kontrast zwischen den Möglichkeiten und der Realität ist also enorm. Es
    gibt unzählige gesellschaftlich relevante Probleme, die an den Hochschulen
    in der Einheit von Forschung und Lehre produktiv bearbeitet werden können,
    und auch den gesellschaftlichen Reichtum (besonders in Hamburg), um das aus
    öffentlicher Hand zu finanzieren. Ob es um den Klimawandel, die
    Massenüberwachung oder das profitgetriebene Gesundheitssystems geht, kann die
    Wissenschaft ein wesentlicher Faktor zur Rationalisierung der Debatte durch
    Aufklärung und der Lösung für eine humane Gesellschaftsentwicklung durch
    Bildung sein. Damit eine in ihrer Schwerpunktsetzung gesellschaftlich
    verantwortliche Universität nicht nur eine Farce ist, braucht es ein steigendes
    Globalbudget für die Hochschulen. Die flächendeckende Ausfinanzierung statt
    einer auf Standortlogik setzende Spitzenförderung ist dafür unabdingbar.

    Was für die Hochschulen gilt, trifft auch auf Schulen, Kitas, sozialen
    Wohnungsbau u.v.m. zu. Der Senat versucht mit seiner Schuldenbremsen-Politik
    den immer offener zutage tretenden Bedarf nach einem Ausbau der sozialen
    Daseinsfürsorge und die soziale Protestbewegung die dafür eintritt
    abzuwiegeln. Dabei bedarf die aktuelle gesellschaftliche Entwicklung aus
    steigender sozialer Ungleichheit und daraus folgender politischer
    Polarisierung einen radikalen Kurswechsel an: Weg von der neoliberalen
    Irrwegen der letzten Jahrzehnte hin zu einer nachhaltigen solidarischen
    Alternative!

    Die Zeit des vermeintlichen "Sparens" muss vorbei sein. Weil wir es ernst
    mit einer demokratischen Entwicklung meinen,  treten wir entschieden dafür
    ein, dass die Fokussierung auf den Standort und damit auf ökonomische
    Partikularinteressen zugunsten einer sozialen Entwicklung für alle überwunden
    wird. Als Studierendenschaft wollen wir diesen Aufbruch vorantreiben, der
    nicht nur die Wissenschaft, sondern alle Bereiche des öffentlichen Lebens
    umfasst.

    Für die weiteren Auseinandersetzungen wird der AStA damit beauftragt, den
    Haushaltsentwurf detaillierter zu sichten und zügig eine Publikation zur
    Kritik an der (Haushalts-)Politik des Senats zu verfassen und in die
    öffentliche Diskussion einzubringen.

    Die Verfasste Studierendenschaft nimmt die Haushaltsverhandlungen zum Anlass,
    die Proteste mit Bündnispartnern in der Stadt und in der Uni zur Lösung der
    Schuldenbremse zu intensivieren. Der AStA wird beauftragt entsprechende
    Aktivitäten vorzubereiten und zu unternehmen.

    Das Studierendenparlament beschließt zur Diskussion und Positionierung zu
    dieser Auseinandersetzung am 6. Dezember eine Vollversammlung einzuberufen.


\end{document}
